\chapter{M005 --- Slab Construction}

\section{Stage 1 --- Mathematical Specification}

\subsection{Scientific Purpose}

Construct a finite slab from a selected surface termination.

\subsection{Transformation}

\[
\Phi_{\mathrm{slab}} :
\mathcal{T}_{\mathrm{surface}}
\rightarrow
\mathcal{S}_{\mathrm{slab}}
\]

\subsection{Mathematical Inputs}

\begin{itemize}
\item Selected Surface Termination
\item Slab thickness specification
\item Boundary-condition specification
\end{itemize}

\subsection{Mathematical Outputs}

\begin{itemize}
\item Slab
\end{itemize}

\subsection{Domain}

Surface-termination space.

\subsection{Codomain}

Slab space.

\subsection{Preconditions}

\begin{itemize}
\item Valid Surface Termination.
\item Positive slab thickness.
\end{itemize}

\subsection{Postconditions}

A finite slab realizing the selected termination.

\subsection{Invariants}

\begin{itemize}
\item Surface termination identity.
\item In-plane lattice.
\item Crystallographic orientation.
\end{itemize}

\subsection{Gauge Freedoms}

Equivalent origin choices that preserve the same slab realization.

\subsection{Notes}

Thickness, stacking, and boundary-condition parameters are part of the
mathematical definition of the slab, not merely implementation details.
